\section{From Reuse Optimization to a Malicious-Reuse Threat Model}
\label{sec:threat-model}

\subsection{Adversary}

Consider an adversary with one or more privileged capabilities:

\begin{itemize}
  \item modifying a compiler optimization or mapping pass;
  \item altering a runtime scheduler's qubit-liveness decisions;
  \item inserting gates into a circuit after user inspection;
  \item modifying pulse-level implementations beneath trusted gate names;
  \item accessing source-preparation controls or classical parameters;
  \item allocating undocumented physical qubits or channels.
\end{itemize}

The adversary may be a compromised dependency, malicious cloud operator,
insider, supply-chain attacker, or high-capability state actor. The threat does
not require that the adversary can violate quantum mechanics. It exploits the
software stack and whatever information is classically accessible or
orthogonally encoded.

\subsection{Attack pattern}

A generic malicious-reuse pattern is:

\begin{enumerate}[leftmargin=1.8em]
  \item \textbf{Find workspace.} Identify qubits that are clean, dead, or soon
  to be reclaimed.
  \item \textbf{Select a secret-dependent signal.} Target basis labels,
  classical parameters, key-dependent branches, syndrome bits, intermediate
  computational-basis values, or source controls.
  \item \textbf{Couple reversibly.} Compute $g(x)$ into hidden workspace using
  CNOT, Toffoli, controlled rotations, or a lower-level pulse substitution.
  \item \textbf{Preserve the checked output.} Avoid touching the victim path,
  uncompute temporary changes, or exploit the fact that the victim checks only
  a restricted measurement basis.
  \item \textbf{Extract.} Measure the attacker register later, route it to an
  unauthorized output, or expose it through timing, crosstalk, or control data.
\end{enumerate}

\begin{cautionbox}[Encryption is not automatically defeated]
Copying ciphertext does not reveal plaintext when the encryption remains
secure. The attack matters when it reaches plaintext, keys, secret preparation
choices, key-dependent intermediates, metadata, or ciphertext that was itself
supposed to remain private. The paper therefore claims a potential covert
\emph{information-flow channel}, not a universal cryptanalytic break.
\end{cautionbox}

\subsection{Why reclaimed qubits are attractive}

Reclaimed qubits are useful to both optimizer and attacker because they are
expected to stop affecting the legitimate computation. An unauthorized pass can
attempt to hide among valid reuse transformations. Moreover, if the user expects
fewer physical qubits after compilation, extra operations on recycled qubits may
not appear as a simple increase in declared width.

This does not make reuse unsafe by default. It means reuse decisions should be
attested and information-flow checked, much as classical compilers and memory
allocators require trust boundaries.


\subsection{From a narrative threat model to a navigable capability map}

The attack pattern above is intentionally broad. To keep broad scope from
becoming vague scope, \cref{sec:capability-matrix} classifies each combination
of attacker access and information target. The overview matrix is hyperlinked to
an evidence registry in \cref{app:capability-registry}. This structure makes it
possible to state three things at once: what an attacker can do directly, what
can be achieved only by changing the objective or access assumptions, and which
cells are contributions demonstrated or proposed by this project.

\subsection{Privilege vocabulary: use, read, retain, export}
\label{sec:privilege-vocab}

Because ``access to a qubit'' spans very different capabilities, the threat
model uses four verbs with precise, non-overlapping meanings.

\begin{description}
  \item[\textbf{Use.}] Apply a unitary or measurement to the qubit within a
  sanctioned protocol step. Use is the only privilege that a well-designed
  protocol needs to grant. It is strictly weaker than read: applying a gate does
  not imply that the classical outcome is retained.

  \item[\textbf{Read.}] Extract a classical value—destructively via measurement
  or indirectly via a correlated classical side-channel—without authorization.
  Read implies that an outcome bit is written to an adversary-accessible
  register. Use does \emph{not} imply read: a controlled-$Z$ in the middle of a
  circuit is a use but produces no classical side effect unless the circuit also
  measures it.

  \item[\textbf{Retain.}] Keep a quantum copy, or a quantum-correlated ancilla,
  beyond the authorized handoff boundary. Retain does not require measuring the
  copy at the time of theft; the adversary may defer measurement to gain
  side-channel advantage later. Read does \emph{not} imply retain: measuring a
  qubit immediately to get a classical bit produces only a classical residue, not
  a quantum copy.

  \item[\textbf{Export.}] Route a retained copy or classical residue to an
  unauthorized classical or quantum output channel. Export is the final step that
  produces exfiltration. An adversary may retain a copy inside the trusted
  boundary indefinitely; export is what moves it outside.
\end{description}

These four privileges are ordered by what they require of the adversary:
$\text{use} \subset \{\text{read, retain}\} \subset \text{export}$. The
principal counterexamples are: (1) a measurement at the end of a legitimate gate
sequence is a use plus a read but not a retain; (2) a coherent ancilla that
accumulates information without any classical measurement is a retain but not a
read until measured. The capability matrix in \cref{sec:capability-matrix} uses
this vocabulary to label each cell's required privilege level.
